\documentclass[10pt,a4paper]{amsart}
\usepackage[margin=19mm]{geometry}
\usepackage{amsmath,amssymb,mathtools}
\usepackage[scr=rsfs,cal=euler]{mathalfa}
\usepackage{xfrac}
\usepackage[unicode,hidelinks,bookmarks=false]{hyperref}
\newcommand{\ie}{i.e.\ }
\newcommand{\cf}{cf.\ }
\newcommand{\resp}{respectively\ }
\newcommand{\Subsection}{\subsection}
\providecommand{\nobreakdash}{\nobreak\hbox{-}}
\setlength{\emergencystretch}{2em}
\multlinegap0pt
\usepackage{bm}
\usepackage{bbm}
\usepackage{dsfont}
\usepackage{xspace}
\usepackage{cancel}
\usepackage{mathtools}
\usepackage{cleveref}
\usepackage{amsthm}
\makeatletter
\@ifundefined{theorem}{\newtheorem{theorem}{Theorem}}{}
\@ifundefined{lemma}{\newtheorem{lemma}[theorem]{Lemma}}{}
\makeatother
\newcommand{\QQ}{\mathbb{Q}}
\newcommand{\eps}{\bm{\varepsilon}}
\newcommand{\mdp}{\mathcal{M}}
\title[On Piecewise Affine Reachability with Bellman Operators]{On Piecewise Affine Reachability with Bellman Operators}
\author{Anton Varonka, Kazuki Watanabe}
\date{Source: arXiv:2502.19923v3 (CC BY 4.0)}
\begin{document}
\maketitle

\noindent\emph{Source and licence.} On Piecewise Affine Reachability with Bellman Operators. Version arXiv:2502.19923v3 (CC BY 4.0), distributed under the Creative Commons Attribution 4.0 International licence, as recorded in the arXiv licence metadata for this exact version and shown on the abstract page https://arxiv.org/abs/2502.19923v3. Full rights notice: licenses/arxiv-2502.19923-CC-BY-4.0.txt.\par\smallskip

\noindent\textit{Editorial scope.} Counted unit: the Lemma whose source label is \texttt{2d-comp-soon} in the article (source0.tex lines 1196--1206, source label \texttt{2d-comp-soon}), delivered together with the article's own complete original proof of it (source0.tex lines 1207--1290, one \texttt{proof} environment of 84 lines). No mathematics is written, repaired or re-derived: statement and proof are the article's own text line for line. Required context retained uncounted: none. Every definition, display and label of those lines is retained unchanged. Omitted from this package and not redistributed: 1--1195, 1291--1562. Nothing inside a retained excerpt is omitted or altered: every retained body is delivered exactly as the article's lines are written, so no omission receipt of any kind accompanies this package. Citations: the retained excerpts carry no citation command at all, so no bibliography entry of the article is transcribed and no reference is redistributed. Reference adaptation: none was needed and none was made; every reference inside the delivered unit resolves inside it, so no unresolved reference or citation is produced. Printed slips kept literal: no printed word, symbol, hypothesis or number of the counted statement or of its proof was altered. Layout and class: the article's own class is \texttt{elsarticle} with packages including amssymb, amsmath, amsthm, cleveref, thm-restate, appendix, mdframed, framed, bm, xspace, microtype, xcolor, tikz, todonotes; the wrapper is the lane's pinned \texttt{amsart} editorial wrapper at 10pt on A4, which restates the theorem-like environments the retained text needs and adds the article's own macro definitions that the retained text needs (\texttt{\textbackslash QQ}, \texttt{\textbackslash eps}, \texttt{\textbackslash mdp}), with extra wrapper packages (bm, bbm, dsfont, xspace, cancel, mathtools, cleveref). The delivered numbering is the unit's own. Assets: no retained span contains a figure environment, a graphics-inclusion command, a PDF-page inclusion, a table or a TikZ picture; no image file is copied into this package, the package declares no asset dependency and no image file is redistributed with it. Licence: version arXiv:2502.19923v3 of this article is distributed under the Creative Commons Attribution 4.0 International licence, as recorded in the arXiv licence metadata for this exact version and shown on the abstract page https://arxiv.org/abs/2502.19923v3. A full rights notice is delivered at licenses/arxiv-2502.19923-CC-BY-4.0.txt. Characters: every retained excerpt is pure ASCII, so the delivered engine, XeLaTeX with its default Latin Modern text font, reports no missing character.
\begin{lemma}\label{2d-comp-soon}
	Given a product family $\{M_{1,1}, \dots, M_{k,\ell}\} \subset \QQ^{2\times 2}$ of substochastic matrices
	and a vector~$\eps = (\varepsilon_1, \varepsilon_2)$. 
	The map~$F$ is defined as above by
	\[F(\bm{v}) = \max_{1\leq i \leq k, 1\leq j \leq \ell} \left(M_{i,j} \cdot\bm{v}\right).\]
	Exactly one of the two statements holds:
	\begin{enumerate}
		\item $F^n(\eps) \neq \bm{0}$ for all~$n$.
		\item $F^2(\eps)$ is comparable with~$\bm{0}$.
	\end{enumerate}
\end{lemma}

\begin{proof}
	We can assume~$\eps\bowtie\bm{0}$, otherwise 2.\ holds trivially.
	Without loss of generality, let $\eps = (\varepsilon_1, \varepsilon_2)$ be such that $\varepsilon_1 < 0$ and $\varepsilon_2 > 0$, i.e., $\eps \in Q_2$.
	The discussion below is driven by the question 
	\begin{center}
		``When are both vectors $F(\eps)$ and $F^2(\eps)$ incomparable with~$\bm{0}$?''.
	\end{center}
	
	Crucially, the components of~$F(\bm{v})$ can be expressed as $\max$ of actions applied:
	\begin{equation}\label{each-act-max}
		F(\bm{v}) = 
		\begin{pmatrix} \max_{1\leq i \leq k}  \alpha_i(\bm{v}) \\ \max_{1\leq j \leq \ell}  \beta_j(\bm{v})
		\end{pmatrix}
	\end{equation}
	
	We first identify one special case: when both a zero $\alpha$- and a zero $\beta$-action exist.
	Then $F(\eps)_1 \geq \bm{0}\cdot\eps = 0$ and $F(\eps)_2 \geq \bm{0}\cdot\eps = 0$,
	thus showing that $F(\eps) \geq \bm{0}$.
	We now proceed under the assumption that zero actions exist for at most one state of~$\mdp$.
	Then $F^{-1}(\bm{0})$ is contained in the union of kernel lines.
	
		Observe that for a singular non-zero matrix~$M$, the vector $M\cdot\eps$ is always comparable with~$\bm{0}$.
		Let $\theta \in [\frac{\pi}{2}, \pi]$ be the angle of the kernel line of~$M_{lo}$.
		In our case distinction, we compare the angles~$\measuredangle \alpha_{lo}$ and $\measuredangle\beta_{lo}$ with $\theta$.
		\begin{enumerate}
			\item $\theta = \max(\measuredangle \alpha_{lo}, \measuredangle \beta_{lo})$: kernel of~$M_{lo}$ is the line with the largest angle. 
			We use~\cref{each-act-max} repetitively: either we have 
			$M_{lo}\cdot\eps \geq \bm{0}$ and hence $F(\eps) \geq \bm{0}$, or $M_{lo} \cdot\eps < \bm{0}$. 
			In the latter case, $F(\eps) < \bm{0}$ 
			because $\alpha(\eps) < 0$ and $\beta(\eps) < 0$ for any non-zero $\alpha,\beta$; 
			and there is a state with no zero actions.
			Therefore, $F(\eps)$ is comparable with~$\bm{0}$.
			\item $\theta < \measuredangle\gamma$ for some action $\gamma \in Act_1^* \cup Act_2^*$: there exists an action whose line has a larger angle than all kernel lines.
			For any $M_{i,j}$ and any non-zero $\bm{x} \in \ker M_{i,j} \cap Q_2$ we have $\gamma(\bm{x}) > 0$. 
			Then either $F(\bm{x})_1 > 0$ or $F(\bm{x})_2 > 0$, and hence $F(\bm{x}) \neq \bm{0}$. 
			This yields $\ker M_{i,j} \cap Q_2 = \{\bm{0}\}$ for all~$i,j$,
			and further, $F^{-1}(\bm{0}) \cap Q_2 = \{\bm{0}\}$. 
			
			We may assume $\gamma$ is an action with the largest angle. 
			Note that its line is \emph{not} a kernel line because $\measuredangle\gamma > \theta$.
			Therefore, this line corresponds to actions of one kind (either $\alpha$-, or $\beta$-). 
			It follows that $\measuredangle\alpha_{lo} \neq \measuredangle\beta_{lo}$ and we proceed with the following case distinction.
			\begin{enumerate}
				\item $\measuredangle \beta_{lo} > \measuredangle \alpha_{lo}$. 
				We show that $F^n(\eps) \not\in Q_4$ for all~$n \geq 0$. 
				We emphasise that no zero $\alpha$-action exists, as this would imply that the line of~$\beta_{lo}$ is a kernel line.
				Now, assume towards a contradiction that~$m$ is the least integer such that $F^m(\eps) \in Q_4$.
				Consider $\bm{y}:= F^{m-1}(\eps)$.
				Since $F(\bm{y})_1 \geq 0$ and $\bm{y} \not\in Q_4$, we have $\alpha_{lo}(\bm{y}) \geq 0$.
				This, in turn, together with $\bm{y} \not\in Q_4$, implies $\beta_{lo}(\bm{y}) > 0$. Therefore, $F(\bm{y})_2 = F^m(\eps)_2 > 0$ contradicting our assumption.
				We have $F^n(\eps) \not\in Q_4$ for all~$n \geq 0$.
				
				\item $\measuredangle \alpha_{lo} > \measuredangle \beta_{lo}$.
				If $\alpha_{lo}(\eps) \leq 0$, then also $\beta_{lo}(\eps) \leq 0$ and so $F(\eps) \leq 0$.
				Moreover, if $\beta_{lo}(\eps) \geq \bm{0}$, then $F(\eps) \geq 0$.
				We now assume $\beta_{lo}(\eps) < 0 < \alpha_{lo}(\eps)$.
				Note that no zero $\beta$-action exists, for the line of $\alpha_{lo}$ is not a kernel line.  %
				Then, $F(\eps) \in Q_4 \setminus \{\bm{0}\}$.
			\end{enumerate}
		\end{enumerate}
		
		From the discussion so far, we know that either $F(\eps)$ is comparable with~$\bm{0}$, 
		or $F^n(\eps) \neq \bm{0}$ for all $n \geq 0$,
		or $F(\eps) \in Q_4\setminus \{\bm{0}\}$ while $F^{-1}(\bm{0}) \cap Q_2 = \{\bm{0}\}$.
		We now investigate the third option.
		Here, we focus our attention on actions~$\alpha_{hi}$ and $\beta_{hi}$ rather than~$\alpha_{lo}$ and $\beta_{lo}$.
		
		
		Let $\eta \in [\frac{\pi}{2}, \pi]$ be the angle of the kernel line of~$M_{hi}$.
		\begin{enumerate}
			\item $\eta = \min(\measuredangle \alpha_{hi}, \measuredangle \beta_{hi})$: kernel of~$M_{hi}$ is the line with the smallest angle. Then 
			either $M_{hi}\cdot F(\eps) \geq \bm{0}$ and hence $F^2(\eps) \geq \bm{0}$, or $M_{hi} \cdot F(\eps) < \bm{0}$. In the latter case, $F^2(\eps) < \bm{0}$ because $\alpha(F(\eps)) < 0$ and $\beta(F(\eps)) < 0$ for any non-zero $\alpha,\beta$. 
			In either case,	$F^2(\eps)$ is comparable with~$\bm{0}$.
			\item $\eta > \measuredangle\gamma$ for some action $\gamma \in Act_1^* \cup Act_2^*$: there exists an action whose line has a smaller angle than all kernel lines.
		For any $M_{i,j}$ and any $\bm{x} \in \ker M_{i,j} \cap Q_4$ we have $\gamma(\bm{x}) > 0$. Then $F(\bm{x}) \neq \bm{0}$, and hence $F^{-1}(\bm{0}) \cap Q_4 = \{\bm{0}\}$.
			
			Recall that we already proved $F^{-1}(\bm{0}) \cap Q_2 = \{\bm{0}\}$ for the case under consideration.
		Moreover, since $F^{-1}(\bm{0})$ is contained in the union of kernel lines and thus in $Q_2 \cup Q_4$, we have $F^{-1}(\bm{0}) = \{\bm{0}\}$.
		Therefore, $F^n(\eps) \neq \bm{0}$ for all~$n\geq 0$.
		\end{enumerate}
		
		In most cases we were able to show $F(\eps)$ is comparable with~$\bm{0}$. This implies $F^2(\eps)$ is comparable with~$\bm{0}$, too.
		In all other cases we either directly showed that $F^2(\eps)$ is comparable with~$\bm{0}$, or derived $F^n(\eps) \neq \bm{0}$ for all~$n$.
\end{proof}

\end{document}
